\chapter{序論}
\label{chap:序論}

\section{背景と目的}
\label{sec:背景と目的}

\subsection{創造的活動におけるAIのパラダイムシフト}
近年，深層学習技術の飛躍的な進歩，とりわけ大規模言語モデル（LLM）や拡散モデルに代表される生成AI（Generative AI）の登場は，人工知能の役割を根本から変えつつある．従来のAIは，認識，分類，最適化といった定型的なタスクを効率化するための「自動化ツール」としての側面が強かった．しかし，現在の生成AIは，文章執筆，画像生成，楽曲制作など，かつては人間の独壇場であった創造的領域（Creativity）において，人間と同等あるいはそれ以上の品質で成果物を生成する能力を獲得している．

この技術的進歩は，人間とAIの関係性において新たなパラダイムシフトをもたらしている．すなわち，人間が抱く明確な完成イメージを効率的に具現化する「道具」としての役割を超え，人間が想起し得なかったアイデアを提示し，創造性を拡張する「パートナー」としての可能性である．

\subsection{一方的な命令から双方向の共創へ}
しかしながら，現状の多くの生成AIシステムにおけるインタラクションは，依然として人間が一方的に命令（プロンプト）を与え，AIがそれを出力するという「主従関係（Master-Slave）」の枠組みに留まっている．この一方通行の構造には，創造的パートナーシップを築く上で2つの本質的な課題が存在する．

第一に，\textbf{「意図の言語化」の壁}である．創造的活動の初期段階において，人のアイデアは往々にして曖昧であり，言語化できない暗黙的なイメージを含んでいる．一方的な命令形式では，ユーザーが自身の意図を完璧に言語化できなければ，望む結果を得ることができない．

第二に，\textbf{「フィードバック」の硬直性}である．AIの出力がユーザーの意図と乖離した場合，プロンプトを書き直して再生成するという不連続な作業が必要となり，これが創造的なフロー状態（Flow State）を阻害する要因となる．

真に創造的な共創を実現するためには，AIが人間の命令を機械的に処理するだけでなく，互いの意図を理解し，刺激し合う「動的な相互作用（Dynamic Interaction）」が必要不可欠である．本研究では，この動的な相互作用を以下の3つの要素を包含するプロセスと定義し，その実現を目指す．

\begin{enumerate}
    \item \textbf{人間の内部状態の変容}: AIとの対話や生成された結果を通じて，ユーザー自身の感じ方や思考，あるいは当初抱いていた目標イメージそのものが変容・深化すること．
    \item \textbf{相互創発性}: ユーザーとAI，互いのアイデアが刺激となり，一人（あるいはAI単体）では到達し得ない新たな創造や表現が生まれること．
    \item \textbf{文脈の共有と記憶}: その場限りの一時的な応答ではなく，過去の対話履歴や修正の経緯（コンテキスト）を踏まえ，互いの意図を深く理解した上での継続的なやり取りが行われること．
\end{enumerate}

\section{音楽生成における課題：可視化の壁と時間性}
\label{sec:why_music}

人とAIが共創を行う際，対象となるドメイン（創作分野）の特性によって，コミュニケーションの様相は大きく異なる．本研究では，数ある創作活動の中でも「音楽（演奏表現）」を対象とする．その理由は，画像生成やテキスト生成と比較した際，音楽が持つ「不可視性」と「時間依存性」という特性が，AIとの意思疎通において特有かつ高度な課題を提起するためである．

\subsection{可視性と一覧性の欠如}
画像生成の領域において，AIの出力結果は常に「可視化」された状態で提示される．ユーザーは生成された画像を視覚的に捉え，全体像（ゲシュタルト）を瞬時に把握することが可能である．また，修正プロセスにおいても，「右上の雲を消して」や「背景の色を明るくして」といった具合に，空間的な位置を指差すように具体的な指示を出すことができる．これは，視覚情報が持つ「一覧性（Synoptic view）」による恩恵である．

対照的に，音楽は本質的に「不可視」のメディアである．波形や譜面として視覚化することは可能であるが，それらはあくまで記号的な表現に過ぎず，そこから想起される「感情」や「ニュアンス」といったクオリアを直接的に目で見ることはできない．ユーザーは，AIが生成した音楽を確認するために，必ず「再生」という時間をかけた行為を必要とする．
この「可視化できない」という性質は，AIへのフィードバックループにおいて重大な障壁となる．ユーザーは「あそこの盛り上がり」や「なんとなく悲しい感じ」といった，物理的に指差し確認が不可能な対象を，言語という不完全な媒体を通してAIに伝え，共有しなければならないからである．

\subsection{時間軸上の不可逆性と文脈依存性}
音楽におけるもう一つの決定的な特徴は，情報が時間軸上に展開される「継時性（Sequentiality）」である．
静止画であれば，画面の端にある要素を変更しても，中心にある被写体の意味合いは大きく変わらない場合が多い．しかし音楽の場合，ある一瞬の音の強弱やタイミングの変更は，その前後のフレーズの聴こえ方，ひいては楽曲全体の物語性（Narrative）に影響を及ぼす．「今の瞬間」の表現は，過去の演奏の文脈の上に成り立っており，未来の展開への期待を含んでいる．

したがって，音楽生成におけるAIとの対話では，単なる「部分的な修正」ではなく，時間的な文脈（Temporal Context）を考慮した「全体的な整合性」の維持が求められる．この時間的な依存関係を，言語を通じてAIと共有し制御することは，静的な画像の生成に比べて遥かに複雑な認知プロセスを要する．

\subsection{感性語の曖昧性と解釈の主観性}
上述した「可視化の困難さ」と「時間的な文脈」の問題は，指示に用いる「言語（プロンプト）」の役割をより一層重いものにする．
画像生成における「赤いリンゴ」という指示は，客観的な色彩定義と結びつきやすく，万人に共通する視覚的イメージが存在する．一方で，音楽における「悲しげに」や「ドラマチックに」といった指示は，テンポ，音色，アーティキュレーションの複合的な組み合わせによって表現され，その正解は受け手個人の感性や文化的背景に強く依存する．
明確な視覚的リファレンスが存在しない中で，いかにしてこの主観的かつ曖昧なイメージをAIとすり合わせるかが，本研究が挑む中心的な問いである．

\section{本研究のアプローチ}
\label{sec:approach}

\subsection{Human-in-the-loopによる感性のすり合わせ}
上述の課題に対し，本研究ではHuman-in-the-loop（人間参加型）システムを用いた，言語を通じた「感性のすり合わせ」プロセスを提案する．
可視化できない音楽表現において，言語はユーザーとAIをつなぐ唯一の共通プロトコルとなる．本システムでは，ユーザーが曖昧な感性語（例：「もっと情熱的に」「ここは繊細に」）を用いてイメージを伝え，AI（大規模言語モデルを活用した解釈モジュール）がそれを具体的な演奏パラメータに変換して提示するアプローチをとる．

\subsection{意味的インタラクションのループ}
本研究の核となるのは，単発の指示生成ではなく，反復的な対話を通じた意味の共有（Grounding）である．
ユーザーはAIが生成した演奏を聴取し，自身のイメージとの差異を再び言語でフィードバックする．このループ構造により，一度の指示では伝達困難な暗黙的なニュアンスを徐々に具体化し，ユーザーの意図に対するAIの理解（文脈の記憶）を深める．

さらに，本システムは単なる指示追従にとどまらず，共有された目標イメージを踏まえた上で，AI側から新たな表現プランを提案する機能を持つ．AIからの予期せぬ提案がユーザーの内部状態に変化を与え，新たな着想を引き出すことで，人とAIの相互創発的な共創を実現する．

\section{本論文の構成}
本論文の構成は以下の通りである．

第\ref{chap:理論・関連研究}章\textbf{「関連研究」}では，創造的AI，音楽情報処理，およびHCIにおけるインタラクションデザインに関する既存研究を概観し，本研究の位置付けを明確にする．特に，音楽生成における感性語の扱いや，共創システムにおける課題について論じる．

第\ref{chap:提案手法}章\textbf{「提案手法」}では，本研究で開発した音楽共創システムの詳細について述べる．MidiCapsを用いた言語から音楽特徴量への変換機構，および対話的な表現生成を実現するアーキテクチャについて詳述する．

第\ref{chap:実験}章\textbf{「評価実験」}では，提案システムの有効性を検証するために行った被験者実験の設計について述べる．プロの音楽家および非専門家を対象とした比較実験の条件設定を示す．

第\ref{chap:結果}章\textbf{「実験結果」}では，実験により得られた定量的データおよび定性的フィードバックを提示する．

第\ref{chap:考察}章\textbf{「考察」}では，実験結果に基づき，提案システムがユーザーの内部状態や創造性に与えた影響について分析する．また，可視化できないメディアにおけるAIとの対話のあり方について議論を深める．

最後に，第\ref{chap:結論}章\textbf{「結論」}において本研究の成果を総括し，今後の展望を述べる．